\subsection{Riemannian Walk (RWalk)}
\label{sssec:rwalk_background}

\gls{rwalk} is a generalisation of \gls{ewc} and \gls{si}: it regularises updates using both local curvature information (Fisher) and trajectory-based importance scores from the optimisation path. This yields an importance-weighted parameter consolidation objective that is updated online during sequential training.

\paragraph{Combined Importance Estimation}

Treat each scalar parameter $\theta_i$ independently (the same update is applied element-wise to tensors). During a task, \gls{rwalk} maintains running statistics $\tilde F_i$ (Fisher) and $\tilde S_i$ (path). Let $g_i=\partial\mathcal{L}/\partial\theta_i$ be the gradient from the minibatch loss used for the current step, $\Delta\theta_i$ be the actual parameter change applied immediately afterwards by the optimiser, and $\alpha\in(0,1)$ and $\varepsilon>0$ govern smoothing and numerical stability.

For curvature, \gls{rwalk} uses an online diagonal Fisher exponential moving average rather than task-end averaging:
%
\begin{equation}\label{eqn:rwalk_fisher_ema}
    \tilde F_i \leftarrow \alpha\, g_i^2 + (1-\alpha)\,\tilde F_i,
    \qquad
    \alpha\in(0,1).
\end{equation}
%
This uses the same diagonal empirical Fisher ingredient $g_i^2$ while exponentially discounting older minibatches, and it couples directly to the Fisher-weighted normaliser in the path update below. By contrast, \gls{ewc} typically consolidates Fisher importance using task-wise averages at task boundaries.

The path statistic accumulates a \emph{loss-change to Riemannian step-size} ratio. Using the first-order loss-change proxy $-g_i\,\Delta\theta_i$ (the per-coordinate analogue of the synaptic contribution in \gls{si}) and a Fisher-weighted squared step $\tfrac{1}{2}\tilde F_i(\Delta\theta_i)^2$, each optimisation step contributes
%
\begin{equation}\label{eqn:rwalk_path_update}
    \tilde S_i \leftarrow
    \tilde S_i +
    \frac{-g_i\,\Delta\theta_i}{\tfrac{1}{2}\tilde F_i(\Delta\theta_i)^2+\varepsilon}
    \, .
\end{equation}
%
Large contributions arise when $-g_i\,\Delta\theta_i$ is large relative to the Fisher-weighted step size, i.e.\ when the update is disproportionately effective at reducing loss compared with its local Riemannian magnitude.

At a task boundary, running values are turned into consolidation weights: the Fisher term is taken as $F_i\leftarrow \tilde F_i$, and the path term is rescaled by a factor $\tfrac{1}{2}$, $S_i\leftarrow \tfrac{1}{2}\tilde S_i$, matching the usual pairing between the $\tfrac{1}{2}$ in the Fisher quadratic form and the discrete path integral. The effective quadratic importance used in~\eqref{eqn:rwalk_loss} is the sum
%
\begin{equation}\label{eqn:rwalk_importance}
    I_i = F_i + S_i
\end{equation}

\paragraph{\gls{cl} Approach}

At task $t$, \gls{rwalk} minimises current-task cross-entropy together with a quadratic penalty around the previous consolidated parameters $\theta_i^*$:
%
\begin{equation}\label{eqn:rwalk_loss}
    \mathcal{L}_\text{RWalk}
    =
    \mathcal{L}_\text{CE}
    +
    \lambda
    \sum_i (F_i + S_i)\,(\theta_i-\theta_i^*)^2
\end{equation}
%
where $\lambda$ controls the stability--plasticity trade-off. Relative to pure Fisher regularisation, the additional path term captures how parameter updates affected optimisation dynamics over the task trajectory, which can improve importance estimates when curvature alone is insufficient.

\paragraph{Inference}

\gls{rwalk} uses the current model directly at inference time. In \gls{til}, task identity is used to select the active task logit slice. In \gls{cil}, prediction is taken from the full logit vector across all seen classes.
